% Standalone Thai abstract, compiled separately from proposal.tex.
% Font: TH Sarabun New (the actual SIPA-released face, in fonts_thsarabunnew/), the KMUTT/Thai-official
% standard. Requires the font installed where fontspec/fc-list can see it (e.g. ~/.fonts/thsarabunnew/).
% Margins match proposal.tex's KMUTT page geometry (top 30mm, left 40mm, right 20mm, bottom 20mm).
% Compile with XeLaTeX (needed for the ICU-based Thai line-breaker, \XeTeXlinebreaklocale):
%   xelatex thai_abstract_standalone.tex
% This machine has no `xelatex` binary preinstalled; a local xetex-based one was built without sudo
% into ~/.local/bin/xelatex (backed by a format file in ~/texmf-var/web2c/xetex/xelatex.fmt).
\documentclass[11pt,a4paper]{article}
\usepackage[a4paper,top=30mm,left=40mm,right=20mm,bottom=20mm]{geometry}
\usepackage{fontspec}
\usepackage{polyglossia}
\setmainlanguage{thai}
\setotherlanguage{english}
\newfontfamily\thaifont{TH Sarabun New}
\setmainfont{TH Sarabun New}
\linespread{1.0}
\XeTeXlinebreaklocale "th"
\XeTeXlinebreakskip=0pt plus 0.1em minus 0.05em
\pagestyle{empty}
\begin{document}
\noindent
\begin{tabular}{@{}p{4.8cm}p{9.7cm}@{}}
หัวข้อโครงงานวิจัย & การถ่ายทอดการเคลื่อนที่ข้ามรูปแบบร่างกายจากแบบจำลองโลกเชิงภาพมุมมองบุคคลที่หนึ่งโดยไม่มีความรู้ล่วงหน้าเกี่ยวกับสัณฐานวิทยา\\[2mm]
หน่วยกิต           & 6\\[2mm]
ผู้ดำเนินโครงงาน   & นางสาวดิษย์ธร สุทธาเวศ\\[2mm]
อาจารย์ที่ปรึกษา   & นายบวรศักดิ์ สกุลเกื้อกูลสุข\\[2mm]
หลักสูตร           & วิศวกรรมศาสตรบัณฑิต\\[2mm]
สาขาวิชา           & วิศวกรรมหุ่นยนต์และระบบอัตโนมัติ\\[2mm]
ปีการศึกษา         & 2569\\
\end{tabular}

\vspace{3mm}
\begin{center}บทคัดย่อ\end{center}
\vspace{1mm}

\noindent ตัวควบคุมการเคลื่อนที่ของหุ่นยนต์ขาที่ฝึกด้วยการเรียนรู้แบบเสริมกำลัง (Reinforcement Learning) จะผูกติดกับ
ร่างกายที่ใช้ฝึกโดยตรง และวิธีการถ่ายทอดข้ามร่างกายที่มีอยู่ในปัจจุบันล้วนต้องได้รับข้อมูลของร่างกายใหม่ก่อนเสมอ ไม่ว่าจะเป็น
ไฟล์ออกแบบ (CAD/URDF) โครงสร้างจลนศาสตร์ ตัวปรับเฉพาะหุ่นยนต์ หรือข้อมูลสาธิตของหุ่นยนต์เป้าหมาย งานวิจัยนี้จึงทดสอบว่า
พฤติกรรมการเคลื่อนที่ที่บันทึกจากหุ่นยนต์ตัวหนึ่งจะขับเคลื่อนหุ่นยนต์อีกตัวที่มีโครงสร้างต่างกันโดยสิ้นเชิงได้หรือไม่ โดยใช้เพียง
วิดีโอจากมุมมองบุคคลที่หนึ่ง (egocentric) ผ่าน \textbf{พิกัดการเคลื่อนที่ของลำตัวร่วม} (ความเร็วไปข้างหน้า ความเร็วด้านข้าง
และอัตราการหมุน ปรับให้ไร้มิติด้วยเลขฟรูด) โดยไม่ใช้แบบจำลองจลนศาสตร์ ไม่ใช้ไฟล์ URDF และไม่ใช้ข้อมูลสาธิตของหุ่นยนต์
เป้าหมาย หุ่นยนต์ที่ใช้คือแมลงกิ่งไม้จำลอง (\emph{Medauroidea extradentata}) ขนาด 18 องศาอิสระ ในโปรแกรม CoppeliaSim
และหุ่นยนต์สี่ขา Unitree B1 ขนาด 12 องศาอิสระ ในโปรแกรม MuJoCo ซึ่งปริภูมิการกระทำของทั้งสองไม่มีมิติหรือความสอดคล้อง
ร่วมกันเลย สถาปัตยกรรมประกอบด้วยตัวเข้ารหัสภาพ V-JEPA2 แบบตรึงค่า (frozen) ร่วมกับแบบจำลองการเปลี่ยนผ่านเชิงผกผัน
(Inverse Transition Model) เชิงไปข้างหน้า (Forward Transition Model) และตัวถอดรหัสการเคลื่อนไหว (Motion Decoder)
ผลการดำเนินงานจนถึงปัจจุบันแสดงว่าพิกัดการเคลื่อนที่ของลำตัวร่วมสามารถถ่ายทอดข้ามร่างกายได้โดยไม่ต้องจับคู่เฟรมหรือปรับ
เป้าหมายระหว่างหุ่นยนต์ทั้งสอง เนื่องจากการเดินมีลักษณะเป็นคาบ มุมมองบุคคลที่สาม (allocentric) จึงทำให้ท่าทางของหุ่นยนต์
บ่งบอกคำสั่งได้เอง แบบจำลองโลกจึงทำนายอนาคตได้โดยใช้การกระทำแฝงที่เรียนรู้มาน้อยมาก ขณะที่มุมมองบุคคลที่หนึ่งทำให้
แบบจำลองกลับมาใช้การกระทำได้จริง ข้อจำกัดที่ยังเหลืออยู่คือแบบจำลองใช้การกระทำได้ในระดับหยาบเท่านั้น ยังไม่สามารถจัด
อันดับความแตกต่างที่ละเอียดหรือทำนายต่อเนื่องในระยะยาวได้

\vspace{2mm}
\noindent คำสำคัญ: การเคลื่อนที่ข้ามรูปแบบร่างกาย / แบบจำลองโลกเชิงภาพ / มุมมองบุคคลที่หนึ่ง / พิกัดการเคลื่อนที่ของลำตัว / V-JEPA2
\end{document}
